\section{Securities Lending in Brief}\label{sec:tour-sbl}

Every stop so far kept one thing fixed: a position is a signed number per unit, and
ownership \emph{is} that number. Securities lending is the one business that breaks the
assumption, because it separates ownership from possession --- a lender still owns
shares that are physically in someone else's hands. This last stop shows the ledger
absorbing that not by adding a new architecture, but by widening one coordinate.

This stop works the same firm's September special dividend --- the \EUR{2} distribution
the option stop already met --- through a book that holds 1{,}000 ACME shares with 400
of them out on loan. A loan is opened before the dividend and recalled after. A scalar
balance can no longer describe the book: it must show that it still owns 1{,}000 shares
for the purpose of profit and loss, while only 600 are available to sell. So the balance
generalises from a scalar to a vector of coordinates, per holder and unit --- one
coordinate for each distinct custody state a share can be in. There are five:

\[
  \vec{w} = (\ \text{own},\ \text{on-loan},\ \text{borrowed},\ \text{collateral-posted},\
  \text{collateral-received}\ ).
\]

Each coordinate passes a physical-action test --- a distinct action moves it and only it
--- and only \emph{own} carries economic value and drives profit and loss. The list is
open at the same seam it is closed: a deployment that permits re-pledging received
collateral adds a \emph{rehypothecated} coordinate the same way, one physical action
moving one coordinate, and this stop does not, so five suffice. Available inventory is
not a stored number at all; it is a projection, computed on read as
$\text{avail} = \text{own} - \text{on-loan} + \text{borrowed}$, so it cannot drift out
of step with the coordinates it is derived from. When ownership and possession
coincide, on-loan and borrowed are zero, avail equals own, and the vector degenerates
back to the scalar of every earlier stop.

Two other events sit at the edges of the loan. A \emph{locate}, the confirmation a
borrower obtains before selling short, enters as a pre-trade check: a stamped
confirmation admitted before the sale, not a balance change. And each loan is its own
unit, carrying its own fee, collateral, and reporting --- so a loan is a first-class
thing in the log, not an annotation on a share balance. The fee is not a figure kept on
the side: it \emph{accrues} over the loan's life, a daily fee on the on-loan
notional booked as periodic cash moves on the loan unit, so the lender's fee income is
itself a fold of recorded moves and reconciles by construction.

\subsection{A loan, opened and recalled}

Opening the loan is one atomic transaction, and the rule that keeps the vector
consistent is the single-coordinate move principle: one atomic move writes exactly one
coordinate of one unit, the \emph{same} coordinate at both of its endpoints --- it
subtracts a quantity from that coordinate at the source wallet and adds the identical
quantity to that coordinate at the destination ($\text{src}[c] \mathrel{-}= q$,
$\text{dst}[c] \mathrel{+}= q$). An operation that must change several coordinates is not
one move but a \emph{transaction}: a list of single-coordinate moves, admitted
all-or-nothing. The open is such a transaction, of three moves.

Two carry the share leg, and each writes its own coordinate against a contra in the
loan's per-relationship virtual wallet. The first marks the book's \emph{on-loan} up by
400 --- $\text{Book.onloan} \mathrel{+}= 400$ against $\text{Virtual.onloan} \mathrel{-}=
400$ --- recording that 400 of the book's shares are now owned-but-lent. The second
records the borrower's possession --- $\text{Borrower.borrowed} \mathrel{+}= 400$ against
$\text{Virtual.borrowed} \mathrel{-}= 400$. The virtual wallet thus carries a contra on
\emph{every} coordinate the transaction writes, which is exactly why the per-unit sum
closes; a lone contra discharging only one of the two coordinates would leave the ledger
400 shares short of balancing. Through both moves the book's \emph{own} is never written
and stays at 1{,}000, and the standing custodian leg that holds the book's shares at
$-1{,}000$ is likewise untouched. The third move is the collateral. Here the collateral
is cash, and cash is value on the \emph{own} coordinate: the move is
$\text{Borrower.own} \mathrel{-}= \EUR{21{,}000}$, $\text{Book.own} \mathrel{+}=
\EUR{21{,}000}$ on the cash unit --- 400 shares marked at the post-split \EUR{50} is
\EUR{20{,}000} of exposure, and a 105\% margin posts \EUR{21{,}000}. It is a single move
on one coordinate, like any cash transfer. The \emph{collateral-posted},
\emph{collateral-received}, and, where a deployment re-pledges, \emph{rehypothecated}
coordinates are named for the non-cash bilateral case, where the collateral is itself
securities to be tracked in custody without being owned; cash collateral needs none of
them, and the return obligation it creates is carried in the loan unit's own state.
After the open:
\[
  \vec{w}_{\text{book}}(\text{ACME}) = (1{,}000,\ 400,\ 0,\ \dots), \qquad
  \vec{w}_{\text{borrower}}(\text{ACME}) = (0,\ 0,\ 400,\ \dots).
\]
The book's available inventory is $1{,}000 - 400 + 0 = 600$; the borrower's is
$0 - 0 + 400 = 400$. Crucially, \emph{no move wrote the book's own in ACME} --- it
stands at 1{,}000 through the whole loan, so the book has not derecognised the shares
and its profit and loss still reads on 1{,}000. That the lender's own is constant across
lending is not a rule enforced after the fact; it holds by construction, because loan
initiation writes on-loan and never own.

Conservation holds here in exactly the sense it held for a plain share, and it is worth
separating the three statements it comprises, because a lent position needs all three.

\emph{Law A --- per-unit balance.} For each unit, the signed sum over \emph{all} slots
that carry it --- the real wallets, the per-relationship virtual wallet, and the
custodian contra --- is zero; equivalently, every transaction nets to zero per unit.
This holds only because the virtual wallet carries a contra on every coordinate a move
writes. Summing the ACME slots the loan touches:
\[
  \underbrace{1{,}000}_{\text{book own}} \; \underbrace{-\,1{,}000}_{\text{custodian}}
  \; + \underbrace{400}_{\text{book on-loan}} \; \underbrace{-\,400}_{\text{virtual on-loan}}
  \; + \underbrace{400}_{\text{borrower borrowed}} \; \underbrace{-\,400}_{\text{virtual borrowed}}
  = 0.
\]
The two virtual contras, one per share coordinate, are what close the sum; drop either
and the line fails to balance by 400.

\emph{Law B --- physical conservation.} The shares physically possessed sum, across all
parties, to the shares in issue --- a constant $N$, here 1{,}000, not zero. Possession is
$\text{own} - \text{on-loan} + \text{borrowed} + \text{collateral-received}$, so a lent
share is counted once, at the borrower who holds it, through the lender's minus: the
book possesses $1{,}000 - 400 = 600$, the borrower $400$, and $600 + 400 = 1{,}000$. (The
collateral-received term is zero here, where the collateral is cash; it carries the
non-cash case, where securities held as collateral are possessed without being owned.)

\emph{Law C --- own is constant.} The lender's \emph{own} is invariant across the loan by
construction, because loan initiation writes \emph{on-loan} and never \emph{own}. This is
the property just seen from the move list, stated as a law: it is why the book's profit
and loss keeps reading on 1{,}000 shares for the whole life of the loan.
\subsection{The special dividend, split by payer}

The loan is open across ACME's September special dividend --- the same \EUR{2} per share
distribution that, in the option stop, moved a written strike. A plain holding takes one
dividend; a lent holding takes two, and the coordinate is what tells them apart. The
book owns 1{,}000 shares, but only 600 are in its possession; the other 400 sit on loan,
held by the borrower on the record date.

The special dividend is a cash event --- no share moves, only cash on the \emph{own}
coordinate --- and it resolves into three cash moves. On the 600 it possesses, the book
receives the \emph{real} dividend, paid by ACME's paying agent:
$600 \times 2 = \EUR{1{,}200}$, carrying a real dividend's tax character --- withholding
and franking as the jurisdiction attaches them to a distribution from the company. On the
400 on loan, ACME pays that same real dividend to whoever holds the shares on the record
date --- the borrower, not the book: $400 \times 2 = \EUR{800}$ from the paying agent to
the borrower. The book is not left short, because the loan obliges the borrower to pass
it a \emph{manufactured} dividend of $400 \times 2 = \EUR{800}$, borrower to lender.

This third leg is a different cash flow in every respect that matters. Its payer is the
borrower, not the issuer; its tax character is a substitute payment rather than a
distribution --- typically carrying no dividend withholding credit, and in many
jurisdictions taxed as ordinary income rather than a qualified dividend in the lender's
hands --- and it is booked on the \emph{on-loan} line, as its own unit and move, precisely
because the on-loan coordinate is what records that these 400 shares are owned-but-lent
and so owe a manufactured, not a real, dividend. Cash conserves across the three moves:
the paying agent is out \EUR{2{,}000}, the book is up \EUR{2{,}000}, and the borrower
nets zero --- it received \EUR{800} from the issuer and passed \EUR{800} on.

The two legs sum to what ownership entitles the book to: $\EUR{1{,}200}$ real plus
$\EUR{800}$ manufactured is $\EUR{2{,}000}$, exactly $\EUR{2} \times 1{,}000$, because
\emph{own} stayed at 1{,}000 through the loan and \emph{own} is what drives profit and
loss. What the loan changed is not how much the book earns but who pays it and in what
character --- and that distinction rides on the coordinate, never smuggled into one
dividend number. A scalar balance could not have told the real \EUR{1{,}200} from the
manufactured \EUR{800}; the vector does, because the on-loan coordinate was there to
carry the difference.

\subsection{Recall, with a deadline}

The recall is not an instantaneous mirror image but a dated obligation, and it generates
no moves at all: it is a notification that flips the loan from \emph{active} to
\emph{recalled} and registers the obligation object the framework already carries. A
recall sets a return deadline; the borrower returning the shares before that deadline is
the discharge, and if the deadline passes with the shares not back, the lender's
\emph{buy-in} --- purchasing them in the market and charging the borrower the cost --- is
the compensation path. A deadline, a discharge predicate, a compensation action, and one
of its terminal states reached by the deadline: this is the obligation object of the
framework, applied to a loan rather than a coupon. The coordinate mirror fires on the
return that settles the recall, not on the recall notice itself.

When the return does complete, it is the mirror of the open, move for move:
$\text{Borrower.borrowed} \mathrel{-}= 400$ against $\text{Virtual.borrowed} \mathrel{+}=
400$; $\text{Book.onloan} \mathrel{-}= 400$ against $\text{Virtual.onloan} \mathrel{+}=
400$; and the \EUR{21{,}000} of cash collateral returns on the \emph{own} coordinate,
$\text{Book.own} \mathrel{-}= \EUR{21{,}000}$, $\text{Borrower.own} \mathrel{+}=
\EUR{21{,}000}$. Afterwards the book's ACME vector is $(1{,}000, 0, 0, \dots)$ and its
available inventory is back to 1{,}000; the borrower's line returns to the zero vector;
the cash collateral is released in full. The available projection moved
$1{,}000 \to 600 \to 1{,}000$ with no stored value that could fall out of step, because
it was a fold of the moves the whole time.

Nothing else changed. The door is the same door, the checks are the same checks, and
the lender's and borrower's differing views of the loan are two folds of the one log,
not two records to reconcile. \emph{Securities lending is accommodated by generalising
the coordinate, not the architecture --- the same door, the same conservation, and views
as folds of the same log carry the position and entitlement mechanics of a loan:
ownership held while possession moves, a dividend split by payer across the lent line,
and a recall discharged or bought in against a deadline --- and they collapse back to the
scalar the moment ownership and possession coincide again.}
